\section{Startup 护城河：速度之后留下什么}

最后一段讨论 startup 护城河。它表面上是创业经验，实际仍然延续前面的“积累问题”。大家常说创业公司优势是快，但尚晏仪指出，快不是完整逻辑。快的价值在于更早积累某些东西，而那些东西才可能成为护城河。

她把 AI 产品的早期竞争放在模型能力拐点里理解。当基础模型突然突破某个临界点，小公司如果最快找到一个场景并推出产品，就能抢到真空期。这个真空期会带来声量、品牌和用户心智。AI 时代的传播成本结构也不同于传统互联网买量：如果一个体验真的突破用户心理防线，用户会自发传播。

但声量并不够。旅行行业的例子说明，产品体验之后会进入供应链竞争。一个 App 做得早、体验好，如果酒店库存不够、价格不稳、履约不强，用户最终还是会回到供应链更强的平台。长期竞争会变成综合战：产品、品牌、供应链、用户行为、组织效率和再投入能力共同形成正循环。

\begin{importantbox}{速度不是护城河，速度带来的资产才是}
快本身不能防御后来者。快只有在转化为品牌、用户习惯、供应链、数据、社区、artifact 或组织能力时，才可能形成长期壁垒。否则，基础模型能力一扩散，功能领先就会被迅速抹平。
\end{importantbox}

谢天宝用 Google 搜索做类比：如果整个互联网在技术上都可以被 index，为什么后来者难以追上 Google？可能因为先做带来了更完整的 index、更丰富的 query 行为、更快的产品迭代和更强的品牌习惯。这个类比提醒我们，先发优势不是玄学，而是一组具体 artifact 的积累。

他们也讨论了什么样的积累更不容易被大厂吃掉。单纯数据可能被更多算力和更多用户追平；但场景 artifact、AI-native 环境、供应链关系、社区信任、组织语义和产品细节更难复制。谷雨提到 NeoCognition 想升级数字环境本身，让环境变得更 AI-native；如果产出就是升级后的环境 artifact，那么早做的积累就是产品本身，而不是间接训练数据。

\begin{knowledgebox}{AI startup 的几类可积累资产}
\begin{tabularx}{\textwidth}{>{\raggedright\arraybackslash}p{0.23\textwidth} Y}
\toprule
资产类型 & 为什么可能形成壁垒 \\
\midrule
品牌与心智 & 用户在新需求出现时先想到你，降低后续获客成本。 \\
供应链与履约 & 在交易型场景中，体验最终要被库存、价格、服务和责任边界兑现。 \\
场景 artifact & 如果产出本身就是可复用环境、workflow 或知识结构，早做会沉淀直接资产。 \\
社区与信任 & 小公司能靠更近的用户关系形成粘性，这种连接不易被纯功能替代。 \\
组织语义库 & 对场景概念、流程和用户需求的长期校准，会提高迭代效率。 \\
\bottomrule
\end{tabularx}
\end{knowledgebox}

结尾处尚晏仪强调人与人的连接。一个产品如果只是工具，用户习惯可能迁移很快；但如果产品意味着用户进入某个社区、确认某种身份、和一群人建立关系，粘性会更强。这是小公司相对于大厂的重要机会：不是在所有资源上取胜，而是在具体人群里建立更密的信任网络。

\subsection{本章小结}

护城河不是一个静态名词，而是动态竞争中的复利结构。AI 初创公司早期的快，必须转化为能积累的资产；中后期的胜负，则取决于产品、供应链、社区、artifact 和组织效率是否能形成正循环。越容易被基础模型平台吸收的能力，越不适合作为长期壁垒；越贴近场景和用户关系的资产，越可能留下。

